\begin{problem}
    Here is a trick question: ``Are there any functions for which the Riemann
    integral converges but the Lebesgue integral diverges?'' 
    Show, however, that the improper Riemann integral
    \[
    \int_0^1 f(x)\,dx
    \]
    of
    \[
    f(x)=
    \begin{cases}
    \dfrac{\pi}{x}\sin\dfrac{\pi}{x}, & \text{if } x\neq 0,\\[1.2ex]
    0, & \text{if } x=0,
    \end{cases}
    \]
    exists and is finite, while the Lebesgue integral is infinite.
    
    \medskip
    
    \noindent
    \textbf{Hint.}
    Integration by parts gives
    \[
    \int_a^1 \frac{\pi}{x}\sin\frac{\pi}{x}\,dx
    =
    \left. x\cos\frac{\pi}{x}\right|_a^1
    -
    \int_a^1 \cos\frac{\pi}{x}\,dx.
    \]
    Why does this converge to a limit as \(a\to 0^+\)?
    
    To check divergence of the Lebesgue integral, consider intervals
    \[
    \left[\frac{1}{k+1},\frac{1}{k}\right].
    \]
    On such an interval the sine of \(\pi/x\) is everywhere positive or everywhere
    negative. The cosine is \(+1\) at one endpoint and \(-1\) at the other. Now use
    the integration by parts formula again and the fact that the harmonic series
    diverges.
\end{problem}

\begin{proof}
    Consider
    \[
    I(a) = \int_a^1 \frac{\pi}{x} \sin\left(\frac{\pi}{x}\right) dx.
    \]
    Let $u = x$, and $dv = \frac{\pi}{x^2}\sin(\frac{\pi}{x})dx$. Using integration by parts, we have
    \[
    \begin{aligned}
    \int_a^1 \frac{\pi}{x} \sin\left(\frac{\pi}{x}\right) dx
    &=
    \left[ x \cos\left(\frac{\pi}{x}\right) \right]_a^1
    -
    \int_a^1 \cos\left(\frac{\pi}{x}\right) dx \\[6pt]
    &=
    \left(1 \cdot \cos(\pi) - a \cos\left(\frac{\pi}{a}\right)\right)
    -
    \int_a^1 \cos\left(\frac{\pi}{x}\right) dx \\[6pt]
    &=
    -1 - a \cos\left(\frac{\pi}{a}\right)
    -
    \int_a^1 \cos\left(\frac{\pi}{x}\right) dx .
    \end{aligned}
    \]
    As $a \to 0^+$, since $\cos\left(\frac{\pi}{a}\right)$ is a bounded function, we get
    \[
    a \cos\left(\frac{\pi}{a}\right) \to 0.
    \]
    Additionally, the improper integral
    \[
    \int_0^1 \cos\left(\frac{\pi}{x}\right)dx
    \]
    converges. Indeed, if \(0<a<b<1\), then
    \[
    \left|\int_a^b \cos\left(\frac{\pi}{x}\right)dx\right|
    \leq \int_a^b 1\,dx
    = b-a.
    \]
    Hence, \(\int_a^1 \cos(\pi/x)\,dx\) is Cauchy as \(a\to 0^+\), and thus it converges to a finite limit. In conclusion, the improper Riemann integral $\lim_{a \to 0^+} I(a)$ exists and is finite.

    To look at Lebesgue integral, consider the sequence of intervals
    \[
    I_k = [\frac{1}{k+1}, \frac{1}{k}].
    \]
    For any $x \in I_k$, we get $\frac{\pi}{x} \in [k\pi, (k+1)\pi]$. Hence, we have
    \[
    \int_{I_k} |f(x)| dx = \left| \int_{\frac{1}{k+1}}^{\frac{1}{k}} \frac{\pi}{x} \sin\left(\frac{\pi}{x}\right) dx \right|.
    \]
    Using integration by parts, we get
    \[
    \begin{aligned}
    \int_{\frac{1}{k+1}}^{\frac{1}{k}} f(x)\,dx
    &=
    \left[ x \cos\left(\frac{\pi}{x}\right) \right]_{\frac{1}{k+1}}^{\frac{1}{k}}
    -
    \int_{\frac{1}{k+1}}^{\frac{1}{k}} \cos\left(\frac{\pi}{x}\right)\,dx \\[6pt]
    &=
    \left(
    \frac{1}{k}\cos(k\pi)
    -
    \frac{1}{k+1}\cos((k+1)\pi)
    \right)
    -
    \int_{\frac{1}{k+1}}^{\frac{1}{k}} \cos\left(\frac{\pi}{x}\right)\,dx \\[6pt]
    &=
    \left(
    \frac{(-1)^k}{k}
    -
    \frac{(-1)^{k+1}}{k+1}
    \right)
    -
    \int_{\frac{1}{k+1}}^{\frac{1}{k}} \cos\left(\frac{\pi}{x}\right)\,dx \\[6pt]
    &=
    (-1)^k
    \left(
    \frac{1}{k}
    +
    \frac{1}{k+1}
    \right)
    -
    \int_{\frac{1}{k+1}}^{\frac{1}{k}} \cos\left(\frac{\pi}{x}\right)\,dx .
    \end{aligned}
    \]
    By triangle inequality, we get
    \[
    \int_{I_k} |f(x)| dx \ge \left( \frac{1}{k} + \frac{1}{k+1} \right) - \int_{\frac{1}{k+1}}^{\frac{1}{k}} \left|\cos\left(\frac{\pi}{x}\right)\right| dx.
    \]
    Since $\left|\cos\left(\frac{\pi}{x}\right)\right| \le 1$, we have
    \[
    \int_{\frac{1}{k+1}}^{\frac{1}{k}} \left|\cos\left(\frac{\pi}{x}\right)\right| dx \le \frac{1}{k} - \frac{1}{k+1} = \frac{1}{k(k+1)}.
    \]
    This implies that
    \[
    \int_{I_k} |f(x)| dx \ge \frac{1}{k} + \frac{1}{k+1} - \frac{1}{k(k+1)} = \frac{2}{k+1}.
    \]
    Hence, we get
    \[
    \int_0^1 |f(x)| dx = \sum_{k=1}^\infty \int_{I_k} |f(x)| dx \ge \sum_{k=1}^\infty \frac{2}{k+1}.
    \]
    This is a harmonic series, which diverges. Hence, we can conclude that
    \[
    \int_0^1 |f(x)| dx = \infty,
    \]
    which implies that the Lebesgue integral is infinite.
\end{proof}

\begin{problem}
    Assume that \(f\) is measurable on \(\mathbb{R}^2\), and assume that at least
    one of the two iterated integrals
    \[
    \int_{\mathbb{R}}
    \left[
    \int_{\mathbb{R}} |f(x,y)|\,dx
    \right]dy
    \qquad\text{or}\qquad
    \int_{\mathbb{R}}
    \left[
    \int_{\mathbb{R}} |f(x,y)|\,dy
    \right]dx
    \]
    exists. Prove the following results:
    \begin{enumerate}
        \item[\textup{a)}] \(f\in L(\mathbb{R}^2)\).
    
        \item[\textup{b)}]
        \[
        \iint_{\mathbb{R}^2} f
        =
        \int_{\mathbb{R}}
        \left[
        \int_{\mathbb{R}} f(x,y)\,dx
        \right]dy
        =
        \int_{\mathbb{R}}
        \left[
        \int_{\mathbb{R}} f(x,y)\,dy
        \right]dx.
        \]
    \end{enumerate}
\end{problem}
\begin{proof}
    \hfill
    \begin{itemize}
        \item [a)] Now we define \(g(x, y) \coloneqq \vert f(x, y) \vert \), then we can give a claim:
        \[
            \iint _{\mathbb{R} ^2} g(x, y) = \int _{\mathbb{R} } \left( \int _{\mathbb{R} } g(x, y) \, \mathrm{d} x  \right) \, \mathrm{d} y = \int _{\mathbb{R} } \left( \int _{\mathbb{R} } g(x, y) \, \mathrm{d} y  \right) \, \mathrm{d} x.    
        \]
        \textit{proof.} For each \(n \in \mathbb{N}\), we define \(C_n \coloneqq [-n, n]^2\), and define a sequence of function \(\left\{ g_n \right\}_{n=1}^{\infty}  \) by 
        \[
            g_n(x, y) = \begin{dcases}
                \min \left\{ g(x, y), n \right\} , &\text{ if } (x, y) \in C_n ;\\
                0, &\text{ if } (x, y) \notin C_n.
            \end{dcases}
        \]   
        Then note that \(g_n \uparrow g\) on \(\mathbb{R} ^2\) and each \(g_n\) is measurable. Also, note that for each \(n \in \mathbb{N}\), \(g_n\) is absolutely integrable since \(C_n\) is bounded and
        \[
        g_n(x, y) \le n \chi_{C_n}(x, y) \implies \int_{\mathbb{R}
        ^2} g_n \le \int_{\mathbb{R}^2} n \chi_{C_n}(x, y) = n m^*(C_n) = n^3 < \infty.
        \]
        By Fubini's theorem, for almost every \(x \in \mathbb{R}\), the map \(y \mapsto g_n(x, y)\) is absolutely integrable, so
        \[
            h_n(x) \coloneqq \int _{\mathbb{R} } g_n(x, y) \, \mathrm{d} y < \infty \quad \text{for almost every } x \in \mathbb{R}.  
        \]        
        Since \(g(x, \cdot)\) is non-negative and measurable, the integral 
        \[
            h(x) \coloneqq \int _{\mathbb{R} } g(x, y) \, \mathrm{d} y \text{ exists, possibly equal to } \infty.  
        \] 
        Since \(g_n \uparrow g\) and \(g_n\) is measurable for all \(n \in \mathbb{N}\), by monotone convergence theorem,
        \[
            \int _{\mathbb{R} } \sup _n g_n(x, y) \, \mathrm{d} y = \sup _n \int _{\mathbb{R} } g_n(x, y) \, \mathrm{d} y \quad \forall x \in \mathbb{R}.  
        \]
        Since 
        \[
            h(x) = \int _{\mathbb{R} } g(x, y) \, \mathrm{d} y = \int _{\mathbb{R} } \sup _n g_n(x, y) \, \mathrm{d} y \quad \text{and} \quad \int _{\mathbb{R} } g_n(x, y) \, \mathrm{d} y = h_n(x),   
        \]
        we have 
        \[
            h(x) = \sup _n h_n(x) \implies h_n \uparrow h \text{ since } \left\{ h_n(x) \right\} \text{ is increasing for all } x \in \mathbb{R}.
        \]
        Now we can apply monotone convergence theorem to \(\left\{ h_n \right\} \) since each \(h_n\) is absolutely integrable by doing Fubini's theorem on \(g_n\). Thus, we have 
        \[
            \sup _n \int _{\mathbb{R} } h_n(x) \, \mathrm{d} x = \int _{\mathbb{R} } \sup _n h_n(x) \, \mathrm{d} x = \int _{\mathbb{R} } h(x) \, \mathrm{d} x = \int _{\mathbb{R} } \left( \int _{\mathbb{R} } g(x, y) \, \mathrm{d} y  \right) \, \mathrm{d} x.     
        \]  
        Now since 
        \[
            \int _{\mathbb{R} } h_n(x) \, \mathrm{d} x = \int _{\mathbb{R} } \left( \int _{\mathbb{R} } g_n(x, y) \, \mathrm{d} y  \right) \, \mathrm{d} x \underset{\text{[Fubini]}}{=} \iint _{\mathbb{R} ^2} g_n(x, y),   
        \]
        we have 
        \[
            \sup _n \int _{\mathbb{R} } h_n(x) \, \mathrm{d} x = \sup _n \iint _{\mathbb{R} ^2} g_n(x, y) = \iint _{\mathbb{R} ^2} \sup _n g_n(x, y) = \iint _{\mathbb{R} ^2} g(x, y)  
        \]
        by monotone convergence theorem again. Thus, we have 
        \[
            \iint_{\mathbb{R} ^2} g(x, y) = \sup _n \int _{\mathbb{R} } h_n(x) \, \mathrm{d} x = \int _{\mathbb{R} } \left( \int _{\mathbb{R} } g(x, y) \, \mathrm{d} y  \right) \, \mathrm{d} x.   
        \]
        Now interchange the role of \(x, y\) and do similar proof, we  have 
        \[
            \iint_{\mathbb{R} ^2} g(x, y) = \int _{\mathbb{R} } \left( \int _{\mathbb{R} } g(x, y) \, \mathrm{d} x  \right) \, \mathrm{d} y.   
        \] 
        Hence, the claim is proved. 

        Now since we know one of the two iterated integrals 
        \[
            \int_{\mathbb{R}}
            \left[
            \int_{\mathbb{R}} |f(x,y)|\,dx
            \right]dy
            \qquad\text{or}\qquad
            \int_{\mathbb{R}}
            \left[
            \int_{\mathbb{R}} |f(x,y)|\,dy
            \right]dx
        \]
        exists, and by the claim:
        \[
            \iint _{\mathbb{R} ^2} \vert f(x, y) \vert  = \int _{\mathbb{R} } \left( \int _{\mathbb{R} } \vert f(x, y) \vert  \, \mathrm{d} x  \right) \, \mathrm{d} y = \int _{\mathbb{R} } \left( \int _{\mathbb{R} } \vert f(x, y) \vert \, \mathrm{d} y  \right) \, \mathrm{d} x,   
        \]
        we can conclude that 
        \[
            \iint_{\mathbb{R} ^2} \vert f(x, y) \vert  
        \]
        exists, so \(f \in L(\mathbb{R} ^2)\). 
        \item [b)] Now since we have shown \(f\) is absolutely integrable, we can use Fubini's theorem to deduce 
        \[
            \iint_{\mathbb{R}^2} f
            =
            \int_{\mathbb{R}}
            \left[
            \int_{\mathbb{R}} f(x,y)\,dx
            \right]dy
            =
            \int_{\mathbb{R}}
            \left[
            \int_{\mathbb{R}} f(x,y)\,dy
            \right]dx.
        \]
    \end{itemize}
\end{proof}

\begin{problem}
    Let
    \[
    f(x,y)=e^{-xy}\sin x\sin y
    \]
    if \(x\ge 0\), \(y\ge 0\), and let \(f(x,y)=0\) otherwise. Prove that both
    iterated integrals
    \[
    \int_{\mathbb R}
    \left[
    \int_{\mathbb R} f(x,y)\,dx
    \right]dy
    \qquad\text{and}\qquad
    \int_{\mathbb R}
    \left[
    \int_{\mathbb R} f(x,y)\,dy
    \right]dx
    \]
    exist and are equal, but that the double integral of \(f\) over \(\mathbb R^2\)
    does not exist. Also, explain why this does not contradict the result in previous exercise.
\end{problem}
\begin{proof}
    We first show that both iterated integrals exist and are equal. First, 
    \[
        \int _{\mathbb{R} } \left[ \int _{\mathbb{R} } f(x, y) \, \mathrm{d} x  \right] \, \mathrm{d} y = \int _0^{\infty} \int _0^{\infty} e^{-xy} \sin x \sin y \, \mathrm{d} x \, \mathrm{d} y = \int _0^{\infty} \sin y \left( \int _0^{\infty} e^{-yx} \sin x \, \mathrm{d} x  \right) \, \mathrm{d} y.      
    \]
    Now fix \(y \ge 0\), define 
    \[
        I_y \coloneqq \int _0^{\infty} e^{-yx} \sin x \, \mathrm{d} x, 
    \] 
    then 
    \begin{align*}
        I_y &\coloneqq \int _0^{\infty} e^{-yx} \sin x \, \mathrm{d} x = \left[ e^{-yx} (-\cos x) \right]_0^{\infty } - \int _0^{\infty} (-y) e^{-yx} (-\cos x) \, \mathrm{d} x \\
        &= (0 - (-1)) - y \int _0^{\infty} e^{-yx} \cos x \, \mathrm{d} x \\
        &= 1 - y \left( \left[ e^{-yx} \sin x \right]_0^{\infty} - \int _0^{\infty} (-y) e^{-yx} \sin x \, \mathrm{d} x   \right) \\
        &= 1 - y \left( (0 - 0) + y \int _0^{\infty} e^{-yx} \sin x \, \mathrm{d} x  \right) = 1 - y^2 I_y.     
    \end{align*}
    This gives
    \[
        I_y = \frac{1}{1 + y^2}.
    \]
    Hence, 
    \[
        \int _0^{\infty} \sin y \left( \int _0^{\infty} e^{-yx} \sin x \, \mathrm{d} x  \right) \, \mathrm{d} y = \int _0^{\infty} \frac{\sin y}{1 + y^2} \, \mathrm{d} y < \int _0^{\infty} \frac{1}{1 + y^2} \, \mathrm{d} y < \int _0^{\infty} \frac{1}{y^2} \, \mathrm{d} y < \infty.      
    \]
    Since \(f\) is symmetric w.r.t. \(x, y\), we can similarly show 
    \[
        \int _{\mathbb{R} } \left[ \int _{\mathbb{R} } f(x, y) \, \mathrm{d} y  \right] \, \mathrm{d} x = \int _0^{\infty} \frac{\sin x}{1 + x^2} \, \mathrm{d} x = \int _0^{\infty} \frac{\sin y}{1 + y^2} \, \mathrm{d} y = \int _{\mathbb{R} } \left[ \int _{\mathbb{R} } f(x, y) \, \mathrm{d} x  \right] \, \mathrm{d} y. 
    \]  
    Now we show the double integral of \(f\) over \(\mathbb{R} ^2\) actually exists, so the problem is wrong. 
    
    If we can show 
    \[
        \int _{\mathbb{R} } \left[ \int _{\mathbb{R} } \vert f(x, y) \vert  \, \mathrm{d} x  \right] \, \mathrm{d} y = \int _0^{\infty} \int _0^{\infty} \vert e^{-xy} \vert \vert \sin x \vert \vert \sin y \vert \, \mathrm{d} x \, \mathrm{d} y < \infty ,  
    \]
    by problem 2 we have 
    \[
        \iint_{\mathbb{R} ^2} f(x, y) = \int _{\mathbb{R} } \left[ \int _{\mathbb{R} } f(x, y) \, \mathrm{d} y  \right] \, \mathrm{d} x = \int _{\mathbb{R} } \left[ \int _{\mathbb{R} } f(x, y) \, \mathrm{d} x  \right] \, \mathrm{d} y,
    \]
    and we have shown 
    \[
        \int _{\mathbb{R} } \left[ \int _{\mathbb{R} } f(x, y) \, \mathrm{d} y  \right] \, \mathrm{d} x = \int _{\mathbb{R} } \left[ \int _{\mathbb{R} } f(x, y) \, \mathrm{d} x  \right] \, \mathrm{d} y < \infty,
    \]
    so this implies the double integral of \(f\) over \(\mathbb{R} ^2\) actually exists.
    
    Now since 
    \begin{align*}
        \int _0^{\infty} \int _0^{\infty} \vert e^{-xy} \vert \vert \sin x \vert \vert \sin y \vert \, \mathrm{d} x \, \mathrm{d} y
    \end{align*}
    can be separated into four parts on \([0, 1] \times [0, 1]\), \([0, 1] \times [1, \infty)\), \([1, \infty) \times [0, 1]\), and \([1, \infty) \times [1, \infty)\), and if we can show all the 4 parts are finite, then we can show
    \[
        \int _0^{\infty} \int _0^{\infty} \vert e^{-xy} \vert \vert \sin x \vert \vert \sin y \vert \, \mathrm{d} x \, \mathrm{d} y < \infty.
    \] 
    \begin{itemize}
        \item Case 1: \([0, 1] \times [0, 1]\), since this region is bounded, so 
        \[
            \int _0^{1} \int _0^{1} \vert e^{-xy} \vert \vert \sin x \vert \vert \sin y \vert \, \mathrm{d} x \, \mathrm{d} y < \infty.
        \]
        \item Case 2: \([0, 1] \times [1, \infty)\), since for \(y \in [0, 1]\), we have \(\vert \sin y \vert \le y \), so we have 
        \[
            e^{-xy} \vert \sin x \vert \vert \sin y \vert \le y e^{-xy} \quad \text{ on } [0, 1] \times [1, \infty).   
        \]
        Hence,
        \begin{align*}
            \int _0^{1} \int _1^{\infty} \vert e^{-xy} \vert \vert \sin x \vert \vert \sin y \vert \, \mathrm{d} x \, \mathrm{d} y &\le \int _0^{1} \int _1^{\infty} y e^{-xy} \, \mathrm{d} x \, \mathrm{d} y = \int _0^1 \left[ - e^{-xy} \right]_1^{\infty} \, \mathrm{d} y \\
            &= \int _0^1 \left[ 0 - (-e^{-y}) \right] \, \mathrm{d} y < \infty.    
        \end{align*}
        \item Case 3: \([1, \infty) \times [0, 1]\), similar as Case 2, just interchange the role of \(x, y\). 
        \item Case 4: \([1, \infty) \times [1, \infty)\).  Note that for \((x, y) \in [1, \infty) \times [1, \infty)\), we have 
        \[
            (x - 1)(y - 1) \ge 0 \iff xy \ge x + y - 1 \ge \frac{x + y}{2}.
        \]
        Thus,
        \[
            e^{-xy} \vert \sin x \vert \vert \sin y \vert \le e^{- \frac{x+y}{2}}.  
        \]
        Now since 
        \begin{align*}
            \int _1^{\infty} \int _1^{\infty} \vert e^{-xy} \vert \vert \sin x \vert \vert \sin y \vert \, \mathrm{d} x \, \mathrm{d} y &\le \int _1^{\infty } \int _1^{\infty} e^{- \frac{x+y}{2}} \, \mathrm{d} x \, \mathrm{d} y = \int _1^{\infty} e^{-\frac{y}{2}} \left( \int _1^{\infty} e^{-\frac{x}{2}} \, \mathrm{d} x  \right) \, \mathrm{d} y \\
            &= \int _1^{\infty} e^{-\frac{y}{2}} \left( \left[ -2 e^{-\frac{x}{2}} \right]_1^{\infty}  \right) \, \mathrm{d} y \\
            &= \int _1^{\infty} e^{-\frac{y}{2}} \left( 0 - \left( -2e^{-\frac{1}{2}} \right)  \right) \, \mathrm{d} y = \left( 2 e^{\frac{1}{2}} \right)^2 < \infty.       
        \end{align*}
    \end{itemize}
    Hence, we have shown that the integral of \(f\) on all 4 regions are finite, and we're done. 
\end{proof}

\begin{problem}
    Suppose that \(\mathcal E\subset C^0([a,b], \R )$ (
     the set of continuous functions $[a,b] \to \R$) is equicontinuous and bounded.
    \begin{enumerate}
        \item[\textup{(a)}] Prove that
        \[
        \sup\{f(x):f\in\mathcal E\}
        \]
        is a continuous function of \(x\).
    
        \item[\textup{(b)}] Show that \textup{(a)} fails without equicontinuity.
    
        \item[\textup{(c)}] Show that this continuous-sup property does not imply
        equicontinuity.
    
        \item[\textup{(d)}] Assume that the continuous-sup property is true for each
        subset \(\mathcal F\subset\mathcal E\). Is \(\mathcal E\) equicontinuous?
        Give a proof or counterexample.
    \end{enumerate}
\end{problem}

\begin{proof}[(a)]
    Suppose that 
    \[
        g(x) = \sup\{f(x):f\in\mathcal E\}
    \]
    for convenience.

    Given $\varepsilon > 0$, since $\mathcal E$ is equicontinuous, there exists $\delta > 0$ such that 
    \[
        \vert x - y\vert < \delta \implies \vert f(x) - f(y) \vert < \varepsilon \quad \text{for all } f \in \mathcal E. 
    \]

    And we claim that 
    \[
        \vert x - y \vert < \delta\implies \vert g(x) - g(y)\vert < \varepsilon 
    \]

    First, we show that $g(x) \le g(y) + \varepsilon$. 

    Since $f(x) \le f(y) + \varepsilon$ and $f(y) \le \sup_{h \in \mathcal E} h(y) = g(y)$, we can get $f(x) \le g(y) + \varepsilon$. Since this inequality hold for all $f \in \mathcal E$, $\sup_{f\in \mathcal E} f(x) \le g(y) + \varepsilon$, and hence $g(x) \le g(y) + \varepsilon$.

    By exchanging the role of $x, y$ we can deduce that $g(y) \le g(x) + \varepsilon$. By combining two inequality, we can conclude that $\vert g(x) - g(y) \vert \le \varepsilon$.
\end{proof}

\begin{proof}[(b)]
    Suppose that $\mathcal E = \{f_n(x) = x^{1/n} : n \in \mathbb{N}\}$ on $[0,1]$. For fixed $n$, $f_n$ is continuous on $[0,1]$, and since $0 \le f_n(x) \le 1$ on $[0,1]$, $f_n$ is bounded. So $\mathcal E\subset C^0([a,b], \R )$.

    However, 
    \[
        g(x)=
        \begin{cases}
        1, & 0 < x \le 1,\\
        0, & x = 0,
        \end{cases}
    \]
    so $g$ is not continuous on $[0,1]$.
\end{proof}

\begin{proof}[(c)]
    Suppose that $\mathcal E = \{f_n(x) = x^{n} : n \in \mathbb{N}\}$ on $[0,1]$. For fixed $n$, $f_n$ is continuous on $[0,1]$, and since $0 \le f_n(x) \le 1$ on $[0,1]$, $f_n$ is bounded. So $\mathcal E\subset C^0([a,b], \R )$.

    Notice that for all $n > m$, $f_n(x) \le f_m(x)$ for all $x \in [0,1]$. So $g(x) = x$, and clearly this is a continuous function.

    However for fixed $\varepsilon > 0$, for all $0 < \delta < 1$, we can pick $N = \lceil log_{(1-\delta)}(1-\varepsilon) \rceil + 1$, such that
    \[
        f_N(1) - f_N(1-\delta) = 1 - (1-\delta)^{N}  > 1 - (1-\delta) ^ {log_{(1-\delta)}(1-\varepsilon)} \ge \varepsilon.
    \]
    And hence $\mathcal E$ is not equicontinuous.

\end{proof}

\begin{proof}[(d)]
    No, we disprove this by providing a counterexample.

    Suppose that $\mathcal E = \{f_n(x) = x^{n} : n \in \mathbb{N}\}$ on $[0,1]$. For fixed $n$, $f_n$ is continuous on $[0,1]$, and since $0 \le f_n(x) \le 1$ on $[0,1]$, $f_n$ is bounded. So $\mathcal E\subset C^0([a,b], \R )$.

    For fixed $\mathcal F \subseteq \mathcal E$ which $\mathcal F$ is non-empty, $\mathcal F = \{x^n : n \in A\}$ where $A \subseteq \mathbb{N}$. Let $q = \min A$.

    Notice that for all $n > m \ge q$, $f_n(x) \le f_m(x)$ for all $x \in [0,1]$. So $g(x) = x^q$, and clearly this is a continuous function.

    However we have proven in (c) that $\mathcal E$ is not equicontinuous.
    

\end{proof}

\begin{problem}
    Is the sequence of functions \(f_n:\mathbb R\to\mathbb R\) defined by
    \[
    f_n(x)
    =
    \cos(n+x)
    +
    \log\left(
    1+\frac{1}{\sqrt{n+2}}\sin^2(n^n x)
    \right)
    \]
    equicontinuous? Prove or disprove.
\end{problem}

\begin{proof}
    We want to show that for all $x, y \in \mathbb{R}$, for all $\varepsilon > 0$, there exists $\delta > 0$ such that
    \[
        \lvert x - y\rvert < \delta \implies \lvert f_n(x) - f_n(y) \rvert \quad \text{for all } n \in \mathbb{N}.
    \]
    Let 
    \[
        h_n(x) = \log\left(
        1+\frac{1}{\sqrt{n+2}}\sin^2(n^n x)
        \right).
    \]

    \begin{claim}
    For fixed $n$, $h_n$ is uniformly continuous over $\mathbb{R}$.
    \end{claim}
    \begin{explanation}
    Suppose that 
    \[
        a_n = \frac{1}{\sqrt{n+2}}, \quad c_n = n^n.
    \]
    Then $h_n(x) = \log\left(1+a_n\sin^2(c_n x)\right)$.
    Then 
    \[
        h'_n(x) = \frac{a_n 2 \sin(c_nx)\cos(c_nx)c_n}{1+a_n\sin^2(c_n x)}.
    \]
    Since $1+a_n\sin^2(c_n x) \ge 1$ and $\vert 2 \sin(c_nx)\cos(c_nx) \vert = \lvert \sin(2c_nx)\rvert \le 1$, 
    \[
        \vert h'_n(x) \vert  = \lvert \frac{a_n 2 \sin(c_nx)\cos(c_nx)c_n}{1+a_n\sin^2(c_n x)} \rvert  \le \frac{a_n \cdot 1 \cdot c_n}{1} = a_n c_n.
    \]
    Then by Mean Value Theorem,
    \[
        \lvert h_n(x) - h_n(y) \rvert \le a_n c_n \lvert x-y \rvert,
    \]
    so $h_n$ is Lipchitz continuous, and hence it is uniformly continuous.
    \end{explanation}

    
    Notice that 
    \[
        \lvert f_n(x) - f_n(y) \rvert \le \lvert cos(n+x) - cos(n+y)\rvert + \lvert h_n(x) - h_n(y)\rvert
    \]

    For fixed $\varepsilon > 0$. Since $\cos(\cdot)$ is 1-Lipchitz continuous function, so $\cos(\cdot)$ is uniformly continuous function. And hence we can pick $\delta_{cos} = \frac{\varepsilon}{2}$ such that 
    \[
        \lvert x - y\rvert \le \delta_{cos} \implies \lvert \cos(x) - \cos(y) \rvert \le \frac{\varepsilon}{2}.
    \]

    And notice that 
    \[
        0 \le h_n(x) = \log\left(
        1+\frac{1}{\sqrt{n+2}}\sin^2(n^n x)
        \right) \le \log\left(
        1+\frac{1}{\sqrt{n+2}}
        \right), \quad \text{which is independent of  } x
    \]
    So
    \[
        h_n \to 0 \text{ uniformly.}
    \]
    And hence we can find some $N$ such that for all $n \ge N$, $\|h_n\|_{\infty} \le \frac{\varepsilon}{4}$, and thus for all $x, y \in \mathbb{R}$ and $n \ge N$,
    \[
        \lvert h_n(x) - h_n(y) \rvert \le \lvert h_n(x) \rvert + \lvert h_y(x) \rvert \le 
        \frac{\varepsilon}{2}.
    \]

    And from previous claim, since $h_i$ is uniformly continuous for all $i \in \mathbb{N}$, for each $i$, we can find $\delta_i$ such that
    \[
        \lvert x -y \rvert \le \delta_i \implies \lvert h_i(x) - h_i(y)\rvert \le \frac{\varepsilon}{2}
    \]

    Then, we take $\delta' = \min\{\delta_1, \delta_2,\cdots, \delta_N, \delta_{cos}\}$ such that
    
    \[
        \begin{aligned}
            \lvert x-y \rvert \le \delta' \implies \lvert f_n(x) - f_n(y) \rvert &\le \lvert cos(n+x) - cos(n+y)\rvert + \lvert h_n(x) - h_n(y)\rvert \\
            &\le \frac{\varepsilon}{2} + \frac{\varepsilon}{2} = \varepsilon. \quad \text{for all } n\in \mathbb{N}. 
        \end{aligned}
    \]

    And hence the sequence of functions $f_n$ is equicontinuous.
    
    
\end{proof}

\begin{problem}
    If \(f:\mathbb R\to\mathbb R\) is continuous and the sequence
    \[
    f_n(x)=f(nx)
    \]
    is equicontinuous, what can be said about \(f\)?
\end{problem}

\begin{proof}
    By the definition of equicontinuity, for every $\epsilon > 0$, there exists $\delta > 0$ such that for all $x, y \in \mathbb{R}$ with $|x - y| < \delta$, we have
    \[
    |f_n(x) - f_n(y)| < \epsilon \quad \forall n \in \mathbb{N}.
    \]
    Substituting $f_n(x)=f(nx)$, we also have
    \[
    |f(nx) - f(ny)| < \epsilon.
    \]
    We now choose two arbitrary points $u, v \in \mathbb{R}$, and we can find a sufficiently large $n$ such that $\frac{|u - v|}{n} < \delta$. Let $x = \frac{u}{n}$ and $y = \frac{v}{n}$, and we have $|x - y| = \frac{|u - v|}{n} < \delta$. By equicontinuity condition, we get
    \[
    \left| f\left(n \cdot \frac{u}{n}\right) - f\left(n \cdot \frac{v}{n}\right) \right| = |f(u) - f(v)| < \epsilon.
    \]
    Since $\epsilon > 0$ was chosen arbitrarily, this holds for any $\epsilon > 0$. Hence, we get
    \[
    |f(u) - f(v)| = 0 \implies f(u) = f(v).
    \]
    Since $u$ and $v$ were also chosen arbitrarily, we can conclude that $f$ is a constant function.
\end{proof}